\section{Normal Koszul models and supported direct image}
\label{ns:sec}

The principal parts in Section~\ref{sec:normal-locality} have a second
description: they form a polynomial module for the normal derivative.
For example, differentiating the class of \(t^{-1}\) generates every
finite principal part.  To extend this observation to several normal
operators, one must retain their action on the entire complex.  A chain
homotopy between two coordinate operators alone does not specify an
equivalence of complexes with all coordinates marked.  We construct
such an equivalence, identify its supported recipient, and then give
finite mutually inverse maps to a normal direct image.

All sums below are algebraic direct sums, and the number \(r\) of
normal directions is finite.  The first comparison holds over an
arbitrary commutative ring.  The subsequent polynomial homotopies and
Weyl-module comparisons require a commutative \(\mathbb Q\)-algebra.
In this section differential modules are \emph{left} modules over the
explicit relative Weyl algebra.  The normal restriction and direct
image defined below use their displayed degrees; no identification
with a differently normalized geometric adjunction is assumed.

\subsection{Finite cutoffs and the supported complex}

Let \(R\) be a commutative ring, let \(W\) be an \(R\)-module, and
let \(N_1,\ldots,N_r\) be commuting \(R\)-linear endomorphisms.
Give \(W\) the structure of a module over
\(S=R[u_1,\ldots,u_r]\) by \(u_iw=N_iw\).
The divided-power algebra \(\Gamma_R(\boldsymbol\tau)\) has basis
\(p_{\boldsymbol n}\), for \(\boldsymbol n\in\mathbb N^r\), and
operations
\[
 p_{\boldsymbol a}p_{\boldsymbol b}
   =\binom{\boldsymbol a+\boldsymbol b}{\boldsymbol a}
      p_{\boldsymbol a+\boldsymbol b},\qquad
 \delta_i p_{\boldsymbol n}=p_{\boldsymbol n-\boldsymbol e_i}.
\]
The binomial coefficient is a product over the coordinates; a basis
element with a negative index is zero.  Each \(\delta_i\) is a
derivation.  Put
\begin{equation}\label{ns:eq:normal-complex}
 C(W)=\Gamma_R(\boldsymbol\tau)\otimes_R
       \Lambda_R(\epsilon_1,\ldots,\epsilon_r)\otimes_R W,
 \qquad
 d=\sum_i\epsilon_i\wedge(N_i-\delta_i),
 \qquad |\epsilon_i|=1.
\end{equation}
The commuting operators and the exterior signs give \(d^2=0\).
Write \(C_N(W)\) for this complex with \(u_i=N_i\) acting on its
last factor.

For \(I=\{i_1<\cdots<i_q\}\), let
\(W_{u_I}=W[(\prod_{i\in I}u_i)^{-1}]\), with
\(W_{u_\varnothing}=W\), and put
\(\epsilon_I=\epsilon_{i_1}\cdots\epsilon_{i_q}\).
The extended \v{C}ech complex
\(\check C_{\boldsymbol u}(W)\) has the summand
\(W_{u_I}\epsilon_I\) in degree \(q\).
Its differential is the sum of localization maps with left exterior
multiplication by \(\epsilon_i\).  This fixes the signs in the
following comparison.

\begin{theorem}[The finite-cutoff comparison]\label{ns:thm:cech}
There is a natural strict \(S\)-linear quasi-isomorphism
\begin{equation}\label{ns:eq:cech-map}
 F:C_N(W)\longrightarrow\check C_{\boldsymbol u}(W),\qquad
 F(p_{\boldsymbol n}\epsilon_I w)=
 \begin{cases}
  \displaystyle\left(\prod_{i\in I}u_i^{-n_i-1}\right)w\epsilon_I,
       &n_j=0\text{ for every }j\notin I,\\[2pt]
  0,&\text{otherwise}.
 \end{cases}
\end{equation}
No flatness, finite-generation, nilpotence or Noetherian hypothesis
on \(W\) or \(R\) is required.
\end{theorem}

\begin{proof}
First take \(r=1\).  For \(d_0\geq0\), the summands with
divided-power index at most \(d_0\) form a two-term subcomplex
\(C_{\leq d_0}\).  Its cohomology is
\begin{equation}\label{ns:eq:cutoff-cohomology}
 \begin{aligned}
 H^0(C_{\leq d_0})&=\ker N^{d_0+1},&
 w&\longmapsto\sum_{n=0}^{d_0}p_nN^nw,\\
 H^1(C_{\leq d_0})&=W/N^{d_0+1}W,&
 \sum_{n=0}^{d_0}p_nb_n\epsilon
   &\longmapsto\sum_{n=0}^{d_0}N^{d_0-n}b_n.
 \end{aligned}
\end{equation}
Indeed, the cycle equations are \(a_{n+1}=Na_n\) and
\(Na_{d_0}=0\).  For a boundary \(b_n=Na_n-a_{n+1}\), where
\(a_{d_0+1}=0\), the second sum telescopes to \(N^{d_0+1}a_0\).
Conversely, if that sum is \(N^{d_0+1}a_0\), the recurrence
\(a_{n+1}=Na_n-b_n\) ends in \(a_{d_0+1}=0\), and constructs
a primitive.  Every cokernel class is represented by \(p_{d_0}w\epsilon\).

The inclusion from cutoff \(d_0\) to cutoff \(d_0+1\) induces
inclusion on the kernels and multiplication by \(N\) on the
cokernels.  Exactness of filtered colimits therefore gives
\[
 H^0(C(W))=\ker(W\longrightarrow W_u),\qquad
 H^1(C(W))=W_u/\operatorname{im}W.
\]
For the second equality the map on the colimit is
\([w]_{d_0}\mapsto[u^{-d_0-1}w]\).  It is surjective.
If this fraction is the image of \(a\in W\), then
\(N^k(w-N^{d_0+1}a)=0\) for some \(k\).
At stage \(d_0+k\), its representative is
\(N^kw=N^{d_0+k+1}a\), so its cokernel class is zero.
This proves injectivity.

In these coordinates, \(F^0(\sum p_na_n)=a_0\) and
\(F^1(\sum p_nb_n\epsilon)=\sum u^{-n-1}b_n\epsilon\).
Its chain equation is the finite telescoping identity
\[
 \sum_{n\geq0}u^{-n-1}(Na_n-a_{n+1})=a_0
 \quad\text{in }W_u.
\]
It induces exactly the two cohomology isomorphisms just proved.

For finite \(r\), apply this comparison successively in each
direction.  All remaining operators commute with the one-variable
map.  Filter the total complex by the exterior degrees in those
remaining directions.  This filtration is finite, and on each
associated piece the one-variable comparison is a quasi-isomorphism.
Induction proves the assertion.  Left exterior multiplication gives
precisely the signs in \eqref{ns:eq:cech-map}.
\end{proof}

Let \(Z=V(u_1,\ldots,u_r)\subset\operatorname{Spec}S\).
The extended \v{C}ech complex represents the supported localization
object \(R\Gamma_Z(W)\); see \cite{stacks-supported}*{Lemma 47.9.1}.
Here this notation means the right adjoint to the inclusion of the
derived category with cohomology supported on \(Z\).
Over a general ring it need not be the right-derived functor of the
ordinary ideal-torsion functor.  The explicit complex and map above
are the comparisons used in this section.

\subsection{A coherent change of coordinate marking}

Now assume that \(R\) is a \(\mathbb Q\)-algebra.  The same complex
\eqref{ns:eq:normal-complex} has another strict \(S\)-action,
\(u_i=\delta_i\); denote it by \(C_\delta(W)\).
Left contraction \(h_i=\iota_{\epsilon_i}\) satisfies
\begin{equation}\label{ns:eq:marking-homotopy}
 dh_i+h_id=N_i-\delta_i.
\end{equation}
To compare all the markings simultaneously, use the polynomial
differential graded algebra
\[
 \Omega=R[s,ds],\qquad |s|=0,\quad |ds|=1,
 \qquad (ds)^2=0,\quad d_\Omega s=ds.
\]
On \(A=\Omega\otimes_R C(W)\), the tensor differential is
\(d_A(\omega\otimes c)=d_\Omega\omega\otimes c+
(-1)^{|\omega|}\omega\otimes dc\).

\begin{theorem}[The marked polynomial roof]\label{ns:thm:roof}
The degree-zero operators
\begin{equation}\label{ns:eq:path-marking}
 U_i=N_i-s(N_i-\delta_i)-ds\,h_i
\end{equation}
define a strict \(S\)-module structure on \(A\).
The last term uses the Koszul tensor sign: on \(f(s)\otimes c\)
it is \(-ds\,f(s)\otimes h_ic\), and it vanishes on
\(ds\,f(s)\otimes c\).
Evaluation at the two endpoints gives a natural roof of strict
\(S\)-linear quasi-isomorphisms
\begin{equation}\label{ns:eq:roof}
 C_N(W)\xleftarrow{\ e_0\ }A
             \xrightarrow{\ e_1\ }C_\delta(W).
\end{equation}
Thus both marked complexes represent the supported object of
Theorem~\ref{ns:thm:cech} in \(D(S)\).
\end{theorem}

\begin{proof}
Write \(L_i=N_i-\delta_i\).  All \(N_i,\delta_i\) commute with
the contractions, and products of two \(ds\,h_i\) terms vanish.
Hence \([U_i,U_j]=0\).  The differential commutators are
\[
 [d_A,N_i-sL_i]=-ds\,L_i,\qquad
 [d_A,-ds\,h_i]=ds\,(dh_i+h_id)=ds\,L_i.
\]
They cancel.  Evaluation \(s=a,ds=0\) gives the claimed coordinate
action at \(a=0,1\), so both endpoint maps are \(S\)-linear chain maps.

On \(\Omega\), define a degree-minus-one map by
\[
 H_a(f(s))=0,\qquad
 H_a(ds\,f(s))=\int_a^s f(v)\,dv.
\]
Polynomial integration exists over \(R\), and
\(d_\Omega H_a+H_ad_\Omega=1-j e_a\), where \(j\) includes
constants.  Tensoring this contraction with \(C(W)\) proves that
each \(e_a\) is a quasi-isomorphism, without any flatness assumption.
The contraction is used as an \(R\)-linear contraction; it is not
asserted to preserve the \(S\)-action.
The inverse derived equivalence is the same roof read in reverse.
\end{proof}

\subsection{Normal derivatives and finite triangular maps}

Let \(\mathcal D_R\) be the relative polynomial Weyl algebra generated
by \(u_i,\partial_i\), with
\[
 [\partial_i,u_j]=\delta_{ij},\qquad
 [u_i,u_j]=[\partial_i,\partial_j]=0,
\]
and with \(R\) central.  Suppose \(W\) is a left
\(\mathcal D_R\)-module, and write \(N_i=u_i|_W\) and
\(D_i=\partial_i|_W\).  Let \(m_i\) be multiplication by
\(p_{\boldsymbol e_i}\) on the divided-power factor.

\begin{lemma}[The derivative correction]\label{ns:lem:derivative}
On \(C_N(W)\) and \(C_\delta(W)\), the derivatives
\begin{equation}\label{ns:eq:corrected-derivative}
 \partial_i^C=D_i-m_i
\end{equation}
give left \(\mathcal D_R\)-module structures compatible with the
differential.  On \(A\), use the same derivatives, acting trivially
on \(\Omega\), and the coordinates \(U_i\).
Every arrow of \eqref{ns:eq:roof} and the map \(F\) is then
\(\mathcal D_R\)-linear.  On localizations of \(W\), derivatives
act by the quotient rule.
\end{lemma}

\begin{proof}
The identity \([\delta_i,m_j]=\delta_{ij}\) gives
\[
 [d,D_i-m_i]=0,\qquad
 [D_i-m_i,N_j]=[D_i-m_i,\delta_j]=\delta_{ij}.
\]
The derivatives commute with one another and with every \(h_j\).
Consequently \([\partial_i^C,U_j]=\delta_{ij}\) and
\([d_A,\partial_i^C]=0\).  Endpoint evaluation preserves them.

If \(i\in I\), the image under \(F\) of the \(-m_i\) term is
\[
 -(n_i+1)u_i^{-n_i-2}
       \prod_{j\in I\setminus\{i\}}u_j^{-n_j-1}w\epsilon_I.
\]
This is the derivative of the denominator in
\eqref{ns:eq:cech-map}; the \(D_i\) term differentiates \(w\).
If \(i\notin I\), \(F\) kills the raised divided-power term and
there is no \(u_i\)-denominator to differentiate.  This proves the
claim on every summand.  Without the term \(-m_i\), the residual
commutator would be \([d,D_i]=-\epsilon_i\wedge\).
\end{proof}

Define the left \(\mathcal D_R\)-module
\[
 \mathcal B_R=\mathcal D_R\big/\sum_i\mathcal D_Ru_i.
\]
Normal ordering identifies it with \(R[z_1,\ldots,z_r]\), where
\(\partial_i\) multiplies by \(z_i\) and
\(u_i=-\partial/\partial z_i\).  It is also a free right
\(R\)-module.  For a complex \(V\) of \(R\)-modules define the
normal direct image by
\[
 i_+V=\mathcal B_R\otimes_R V.
\]
Define the normal restriction by
\(i^!W=R\operatorname{Hom}_{\mathcal D_R}(\mathcal B_R,W)\).
Its explicit representative is
\begin{equation}\label{ns:eq:normal-restriction}
 K(W,N)=\left(\Lambda_R(\epsilon_1,\ldots,\epsilon_r)
               \otimes_R W,\ \sum_i\epsilon_i\wedge N_i\right),
\end{equation}
in degrees zero through \(r\).

\begin{theorem}[The supported direct-image comparison]
\label{ns:thm:pushforward}
Put \(z^{[\boldsymbol n]}=\prod_i z_i^{n_i}/n_i!\), and use
multiindex inequalities and factorials componentwise.
There is a strict \(\mathcal D_R\)-linear chain isomorphism
\(\Phi:i_+K(W,N)\to C_\delta(W)\), with inverse \(\Psi\), given by
\begin{align}
 \Phi(z^{[\boldsymbol n]}\epsilon_Iw)
  &=\sum_{\boldsymbol j\leq\boldsymbol n}(-1)^{|\boldsymbol j|}
       p_{\boldsymbol j}\epsilon_I
       \frac{D^{\boldsymbol n-\boldsymbol j}}
            {(\boldsymbol n-\boldsymbol j)!}w,
       \label{ns:eq:phi}\\
 \Psi(p_{\boldsymbol n}\epsilon_Iw)
  &=\sum_{\boldsymbol j\leq\boldsymbol n}(-1)^{|\boldsymbol j|}
       z^{[\boldsymbol j]}\epsilon_I
       \frac{D^{\boldsymbol n-\boldsymbol j}}
            {(\boldsymbol n-\boldsymbol j)!}w.
       \label{ns:eq:psi}
\end{align}
Together with \eqref{ns:eq:roof} and \eqref{ns:eq:cech-map}, these
maps give a natural equivalence on the ambient normal space,
\begin{equation}\label{ns:eq:supported-equivalence}
 i_+i^!W\simeq R\Gamma_Z(W)\quad\text{in }D(\mathcal D_R).
\end{equation}
Every sum defining a map on an element is finite, independently of
the derivative action on \(W\).
\end{theorem}

\begin{proof}
Normal ordering makes \(\mathcal D_R\) free as a right
\(R[\boldsymbol u]\)-module, with derivative monomials as basis.
The Koszul resolution of the coordinate quotient over
\(R[\boldsymbol u]\) is exact over every \(R\): multiplication
by the next polynomial variable is injective after the preceding
variables have been set to zero.  Extension by the free right
module \(\mathcal D_R\) gives a finite free left
\(\mathcal D_R\)-resolution of \(\mathcal B_R\); its maps are
right multiplications by the coordinates, with Koszul signs.
Applying \(\operatorname{Hom}_{\mathcal D_R}(-,W)\) gives
\eqref{ns:eq:normal-restriction}.  Freeness of \(\mathcal B_R\)
over \(R\) also justifies the ordinary tensor in \(i_+\).

Write \(L_i=D_i-m_i\).  Formula~\eqref{ns:eq:phi} is equivalently
\[
 \Phi(z^{[\boldsymbol n]}\epsilon_Iw)
       =\frac{L^{\boldsymbol n}}{\boldsymbol n!}(p_0\epsilon_Iw).
\]
Indeed, \(m_i^jp_0=j!p_{j\boldsymbol e_i}\).
This expression intertwines \(\partial_i\) with \(L_i\).
The identity \([\delta_i,L_j]=-\delta_{ij}\) intertwines the
coordinate actions \(-\partial/\partial z_i\) and \(\delta_i\).
The differential on \(p_0\epsilon_Iw\) is the Koszul differential
of \(K(W,N)\), and it commutes with every \(L_i\).
Thus \(\Phi\) is a \(\mathcal D_R\)-linear chain map.

The inverse formulas are finite triangular substitutions
\(z_i=D_i-m_i\) and \(m_i=D_i-z_i\).
More explicitly, in either composite the coefficient indexed by
\(\boldsymbol\ell\leq\boldsymbol n\) contains, for each coordinate,
the factor
\[
 \frac{D_i^{n_i-\ell_i}}{(n_i-\ell_i)!}
 \sum_{j=\ell_i}^{n_i}(-1)^{j+\ell_i}
          \binom{n_i-\ell_i}{j-\ell_i}.
\]
It is zero if \(\ell_i<n_i\), and is one if \(\ell_i=n_i\).
This proves both inverse identities, without an infinite series.
Finally the actual chain of comparisons is
\[
 i_+K(W,N)\xrightarrow{\ \Phi\ }C_\delta(W)
     \xleftarrow{\ e_1\ }A
     \xrightarrow{\ e_0\ }C_N(W)
     \xrightarrow{\ F\ }\check C_{\boldsymbol u}(W).
\]
Each arrow is a strict \(\mathcal D_R\)-linear quasi-isomorphism.
Their inverses in the derived category give the inverse of
\eqref{ns:eq:supported-equivalence}.
\end{proof}

\subsection{Cohomology maps and the first normal pole}

In one direction, put
\(T=\ker(W\to W_u)\) and \(Q=W_u/\operatorname{im}W\).
The preceding comparison identifies both cohomology modules and
their derivative actions.

\begin{proposition}[Explicit cohomology and inverse maps]
\label{ns:prop:cohomology}
There are isomorphisms of left Weyl modules
\begin{align}
 R[z]\otimes_R\ker N&\longrightarrow T,&
 z^n\otimes w&\longmapsto D^nw,\label{ns:eq:torsion-map}\\
 \Theta:R[z]\otimes_R\operatorname{coker}N&\longrightarrow Q,&
 z^n\otimes[w]&\longmapsto D^n[u^{-1}w].\label{ns:eq:quotient-map}
\end{align}
On their sources \(u=-\partial_z\) and \(\partial_u=z\).
Both inverses are finite on each element.
\end{proposition}

\begin{proof}
After forgetting the coordinate markings, the maps \(e_0\) and
\(e_1\) are chain homotopic: integration of the \(ds\)-coefficient
from zero to one gives a homotopy from \(e_0\) to \(e_1\).
They therefore induce the same map on the common underlying
\(R\)-cohomology.  Thus the cohomology maps of the full roof can
be computed by the composite of underlying \(R\)-chain maps
\(F\Phi\).  The strict Weyl-linear comparison remains the roof
in Theorem~\ref{ns:thm:pushforward}.
In degree one the quotient rule expands \(D^n[u^{-1}w]\) into
exactly the triangular sum.
Replacing \(w\) by \(w+Nv\) adds derivatives of the image of
\(v\in W\), so the second formula is well defined.
Theorem~\ref{ns:thm:pushforward} already proves that these maps are
isomorphisms.  We give the inverses explicitly.

For any Weyl module \(V\) on which \(u\) is locally nilpotent, set
\begin{equation}\label{ns:eq:projector}
 E_V(v)=\sum_{j\geq0}\frac{D^ju^jv}{j!}.
\end{equation}
The sum is finite.  From
\(uD^j=D^ju-jD^{j-1}\), one obtains \(uE_V=0\) and
\(E_V|_{\ker u}=1\).  The reconstruction identity is
\begin{equation}\label{ns:eq:finite-reconstruction}
 v=\sum_{n\geq0}\frac{(-1)^n}{n!}D^nE_V(u^nv).
\end{equation}
After expansion, the coefficient of \(D^ku^k\) is
\(\sum_{n=0}^k(-1)^n/(n!(k-n)!)\), which is zero for \(k>0\).
All double sums are finite because \(u^kv=0\) for large \(k\).
For \(a\in\ker u\), the same binomial identity gives
\(E_VD^ma=0\) when \(m>0\).  Applying \(E_Vu^n\) to a finite
sum \(\sum_mD^ma_m\) therefore recovers \((-1)^nn!a_n\).
The decomposition is unique, with inverse
\begin{equation}\label{ns:eq:inverse-reconstruction}
 v\longmapsto\sum_{n\geq0}\frac{(-1)^n}{n!}
       z^n\otimes E_V(u^nv).
\end{equation}

Both \(T\) and \(Q\) have locally nilpotent coordinate action.
For \(T\), \(\ker u=\ker N\), giving the inverse to
\eqref{ns:eq:torsion-map}.  For \(Q\), the map
\[
 \theta_0:\operatorname{coker}N\longrightarrow\ker(u:Q\to Q),
 \qquad [w]\longmapsto[u^{-1}w]
\]
is surjective: if \(uv\) is the image of \(w\in W\), then
\(v=u^{-1}w\) in \(W_u\).  Its kernel is represented by
\(NW+T\).  Formula~\eqref{ns:eq:finite-reconstruction} on \(T\)
gives \(uT=T\), since for \(a\in\ker u\)
\[
 D^na=u\left(-\frac{D^{n+1}a}{n+1}\right).
\]
Hence \(T\subset NW\), and \(\theta_0\) is injective.
The inverse of \(\Theta\) is
\[
 q\longmapsto\sum_{n\geq0}\frac{(-1)^n}{n!}z^n\otimes
              \theta_0^{-1}\bigl(E_Q(u^nq)\bigr).
\]
\end{proof}

Return to the coefficient ring \(F\) and power-series module
\(B=F[[t]]\) of Section~\ref{sec:normal-locality}.
In the preceding results take the base ring to be \(R=F\), take
\(W=B\), and let \(u\) act by multiplication by \(t\) and
\(D=\partial_t\).  These are \(F\)-linear operators with
\([\partial_t,t]=1\); the derivative is not \(B\)-linear.
Multiplication by \(t\) is injective and has cokernel \(F\).
Thus the supported complex \([B\to B[t^{-1}]]\), which is the
augmentation ideal of \(\A(B)\) as a complex, has
\begin{equation}\label{ns:eq:principal-part-weyl}
 H^1\cong F[z],\qquad
 z^n\longmapsto\partial_t^n[t^{-1}]
          =(-1)^nn![t^{-n-1}].
\end{equation}
The polynomial \(1\) represents the first nonzero pole.
This is a polynomial Weyl-module comparison.  It specifies no
augmented \(F[[t]]\)-algebra structure on the polynomial endpoint,
no coefficient-bar map, and no vertex comparison.
In particular it does not supply a primitive for
\(\theta_b/t-\theta_c/t\) in the independent interval carriers of
Theorem~\ref{nl:thm:obstruction}.  The pole retained by
\eqref{ns:eq:principal-part-weyl} is precisely the class used there.

\subsection{A level family and its derived special fibre}

The point-valued restriction and its ambient direct image also
distinguish two specializations of a parameter family.
Take
\begin{equation}\label{ns:eq:level-family}
 R=\mathbb C[\kappa],\qquad W=R[t],\qquad
 N=\kappa t-\partial_t,\qquad D=t.
\end{equation}
Here \([D,N]=1\), so \(u=N\), \(\partial_u=D\) defines a
left relative Weyl module.  Write \(G=\operatorname{coker}N\).

\begin{theorem}[The ambient special fibre]\label{ns:thm:level}
For \eqref{ns:eq:level-family}, the map \(N\) is injective and
\begin{equation}\label{ns:eq:level-cokernel}
 G\cong R[\kappa^{-1}]\oplus R[\kappa^{-1}]/R.
\end{equation}
Under the second summand, \([t]\) corresponds to
\(\kappa^{-1}\bmod R\).  The supported complex has only
degree-one cohomology
\[
 Q=W_u/W\cong R[z]\otimes_R G.
\]
There is an exact sequence
\begin{equation}\label{ns:eq:level-exact}
 0\longrightarrow\mathbb C[z]\xrightarrow{\ \beta\ }Q
       \xrightarrow{\ \kappa\ }Q\longrightarrow0,
 \qquad
 \beta(v)=v(\partial_u)[u^{-1}t].
\end{equation}
Consequently \(Q/\kappa Q=0\), whereas
\(\operatorname{Tor}_1^R(Q,R/(\kappa))\cong\mathbb C[z]\).
The derived special fibre of the supported complex is
\(\mathbb C[z]\) in degree zero.  It is the normal direct image
of the degree-zero point-valued restriction \(\mathbb C\).
\end{theorem}

\begin{proof}
For a nonzero polynomial with leading term \(a_n(\kappa)t^n\),
the leading term of its image under \(N\) is
\(\kappa a_n(\kappa)t^{n+1}\), which is nonzero.
Thus \(N\) and all its powers are injective, and \(W\to W_u\)
is injective.  The relations in \(G\) are
\[
 \kappa[t^{n+1}]=n[t^{n-1}]\quad(n\geq0).
\]
They involve even and odd powers separately.  With
\((-1)!!=0!!=1\), put
\[
 a_m=\frac{[t^{2m}]}{(2m-1)!!},\qquad
 b_m=\frac{[t^{2m+1}]}{(2m)!!}.
\]
Their presentations are \(\kappa a_{m+1}=a_m\), and
\(\kappa b_0=0,\ \kappa b_{m+1}=b_m\), respectively.
The maps \(a_m\mapsto\kappa^{-m}\) and
\(b_m\mapsto\kappa^{-m-1}\bmod R\) identify these presentations
with the two summands of \eqref{ns:eq:level-cokernel}.
Both are \(\kappa\)-divisible, and
\(G[\kappa]=\mathbb C\cdot[t]\), the one-dimensional span of the
class \([t]\).  Applying the exact functor \(R[z]\otimes_R-\)
and Proposition~\ref{ns:prop:cohomology} proves
\eqref{ns:eq:level-exact} and the assertions about kernel and cokernel
of \(\kappa\).  The two-term free resolution of \(R/(\kappa)\)
identifies that kernel with the stated \(\operatorname{Tor}_1\).

The connecting class can also be seen before taking supported
cohomology.  For \(v(t)\in\mathbb C[t]\), its horizontal cycle
at \(\kappa=0\) is \(\sum_{j\geq0}p_j(-\partial_t)^jv\).
Lifting this finite sum to \(R[t]\), applying the differential
and dividing by \(\kappa\) gives
\begin{equation}\label{ns:eq:beta}
 \beta(v)=\left[t\sum_{j\geq0}p_j(-\partial_t)^jv\,\epsilon\right]
       \quad\text{in }H^1(C(W)).
\end{equation}
For \(v=t^n\), its image under \(F\) is
\[
 \sum_{j=0}^n(-1)^j\frac{n!}{(n-j)!}
          u^{-j-1}t^{n-j+1}
       =\partial_u^n(u^{-1}t)\quad\bmod W.
\]
The equality is the quotient rule with \(D=t\).
It identifies \eqref{ns:eq:beta} with the injection in
\eqref{ns:eq:level-exact}, under \(v(t)\leftrightarrow v(z)\).

The two terms of \(C(W)\) are \(R\)-free, so ordinary reduction
of this complex computes its derived reduction.  At \(\kappa=0\),
\(N=-\partial_t\) is locally nilpotent.  The kernel and cokernel
calculation in \eqref{ns:eq:cutoff-cohomology}, followed by the
colimit, gives \(H^0(C(W)/\kappa)=\mathbb C[t]\) and
\(H^1(C(W)/\kappa)=0\).  Under \eqref{ns:eq:torsion-map} this is
\(\mathbb C[z]\), with \(z\) corresponding to multiplication by
\(t\).  On the point, the normal restriction at \(\kappa=0\) is
\[
 [\mathbb C[t]\xrightarrow{-\partial_t}\mathbb C[t]],
\]
whose cohomology is \(\mathbb C\) in degree zero, since
differentiation is surjective and has the constants as kernel.
Its direct image is exactly \(\mathbb C[z]\).
Over \(R\), the supported complex is \(Q[-1]\); the
\(\operatorname{Tor}_1\) term therefore supplies its degree-zero
derived special fibre.  Taking only \(Q/\kappa Q\) loses this class.
\end{proof}

\subsection{When the marking roof preserves all vertex products}

The Weyl-module comparison has used linear operators.  A comparison
of fields requires an additional hypothesis on every product.
Let \(R\) again be a commutative \(\mathbb Q\)-algebra, and let
\(W\) be a differential graded vertex superalgebra over \(R\).
Suppose the commuting degree-zero operators \(N_i\) differentiate
every integer-indexed product, commute with the internal differential
and every translation, and kill the vacuum.  Give
\(\Gamma_R(\boldsymbol\tau)\), \(\Lambda_R(\boldsymbol\epsilon)\)
and \(\Omega\) their regular scalar vertex structures with zero
vertex translations.  Include the internal differential of \(W\)
in \eqref{ns:eq:normal-complex}, with the usual tensor signs.

\begin{proposition}[Every-product compatibility]\label{ns:prop:vertex}
Under these hypotheses, \eqref{ns:eq:roof} is a roof of
quasi-isomorphisms of differential graded vertex superalgebras with
commuting marked derivations.  Its maps preserve every integer-indexed
product, the vacuum and all translations.  No local nilpotence is
needed for this roof.

If each \(N_i\) is locally nilpotent, there is in addition a strict
quasi-isomorphism to either endpoint,
\begin{equation}\label{ns:eq:horizontal}
 I:W\longrightarrow C(W),\qquad
 I(w)=\sum_{\boldsymbol n\in\mathbb N^r}
               p_{\boldsymbol n}N^{\boldsymbol n}w.
\end{equation}
It preserves all the same vertex structure and intertwines the
markings \(N_i\) on \(W\) with either \(N_i\) or \(\delta_i\)
on \(C(W)\).
\end{proposition}

\begin{proof}
Each \(N_i\) and \(\delta_i\) is a derivation of every
scalar-extended product.  The contraction \(h_i\) is an odd
derivation of the exterior scalar algebra.  Multiplying it by the
odd scalar \(ds\), with the tensor sign, makes an even derivation.
Thus every \(U_i\) in \eqref{ns:eq:path-marking} differentiates
every product.  It kills the vacuum and commutes with translations.
The contraction supercommutes with the internal differential, so
the commutator proof of Theorem~\ref{ns:thm:roof} still applies.
Endpoint evaluation is a unital differential graded scalar-algebra
map.  It therefore preserves all vertex products and translations;
its underlying quasi-isomorphism is the polynomial contraction
already proved.

For the second assertion, commutation and local nilpotence give,
for each \(w\), a finite rectangular bound on
\(N^{\boldsymbol n}w\).  Hence \eqref{ns:eq:horizontal} is a
finite sum and satisfies \(dI=I d_W\), by cancellation of
\(N_i\) against the index shift \(\delta_i\).
It intertwines both markings by the same shift.
For every product index \(k\), iterated Leibniz gives
\[
 N^{\boldsymbol n}(a_{(k)}b)
   =\sum_{\boldsymbol p+\boldsymbol q=\boldsymbol n}
        \binom{\boldsymbol n}{\boldsymbol p}
        (N^{\boldsymbol p}a)_{(k)}(N^{\boldsymbol q}b).
\]
The divided-power multiplication gives precisely these coefficients
in \(I(a)_{(k)}I(b)\).  Vacuum and translation compatibility follow
from the corresponding hypotheses on \(N_i\).
For a locally nilpotent coordinate, localization is zero.
Thus the extended \v{C}ech complex is just \(W\), and
\(FI=\id_W\).  The proof of Theorem~\ref{ns:thm:cech} applies
degreewise to the internal complex: its finite exterior filtration
still proves that \(F\) is a quasi-isomorphism.
It follows that \(I\) is a quasi-isomorphism as well.
\end{proof}

This proposition applies to the normal vertex complex with its
stated derivations.  It assigns no vertex structure to a general
\v{C}ech localization or to \(i_+K(W,N)\), and the triangular
map \(\Phi\) is not asserted to be multiplicative.
The supported module equivalence and the marked vertex roof are
therefore separate inputs to a possible collision comparison.
Neither changes the coefficient multiplication or supplies the
missing ternary primitive in Section~\ref{sec:normal-locality}.
